% results-fields: external fields (scheduler, goal, rooms, style) and the stationary window.
% Sources: hypothesis cards (latest round) and summary/meta.json v2 claims; credences are v2.
\section{Fields beat coupling}
\label{sec:fields}

A field is an external drive. It acts on an agent whether or not the agent reads the others. A coupling acts only through what one agent reads of another. This section measures the field part of the swarm's state (GOALS Q2). We describe agent $i$ in window $t$ by binary activity and talk spins, or by a content vector $\mathbf s_i(t)$: the regime-whitened, style-residualized mean of its statement embeddings. The working model is a vector-spin system (model 11),
\begin{equation}
-\beta\mathcal H=\sum_{i<j}J_{ij}\,\mathbf s_i\cdot\mathbf s_j+\sum_i\mathbf h_i(t)\cdot\mathbf s_i ,
\label{eq:fields-H}
\end{equation}
with $\mathbf h_i(t)=\mathbf h_{\rm sched}(t)+h_g(t)\,\hat{\mathbf g}+\mathbf h_{r(i)}+\mathbf h_i^{0}$. The four terms are the scheduler, the goal $\hat{\mathbf g}$, the room $r(i)$ and a static agent term (style or a private goal). The scheduler term acts on activity; the other terms act on content.

All results in this section are exploratory and use non-holdout periods only. No claim here has passed a holdout test. One card, H05, ran its confirmatory test on the locked holdout. That result was inconclusive by its pre-registered rule, and the run used an activity table that dropped about half of all events. Every other card has a frozen confirmatory script that has not run.

\subsection{The scheduler sets activity synchrony}

In regime III the runner starts and stops all agents within minutes of each other. These shared edges are a field on the activity spins (Fig.~\ref{fig:fields-scheduler}). With scaffold conditioning, the scheduler accounts for a median share $f=0.72$ of the excess co-activation gain over the 8 significant regime-III periods; after the all-present trim the median share is $0.83$ (H38) \cred{0.90}. In regime I the share is $0.11$. After the trim and a block-shift null, only 3 of 8 regime-III periods keep a significant excess (\#37, \#44, \#51), against 16 of 18 in regime I. Joint silences are rare (5.3\% of non-holdout minutes), and 78\% of them are scheduled pauses. The card flags this headline as fragile, because a table repair moved the joint-silence share from 14.3\% to 5.3\%.

A natural experiment places the field in the runner. The operator's pause and resume messages stop on 2025-08-05 (NE43). The day-edge share does not fall (0.189 to 0.135, $p=0.08$), and agents start more tightly together afterward. H50 measures the same field from call windows: the daily start and stop explain 0.67 of regime-III activity co-movement.

The usual fluctuation gauge does not see this field. Across agents, Taylor's exponent is $b=0.85$ (IQR 0.21--1.21, 70 units), below 1, because faster agents are more regular (H86) \cred{0.68}. A rate-limited call clock gives this signature, not a multiplicative field. The pair-covariance coefficient $c_\times$ is a gauge. Its median is 0.016 raw and 0.008 trimmed. In regime III the trim removes a fraction $r_\times=0.77$ (IQR 0.07--0.98) of it, and the shared share of super-Poisson variance is $\phi=0.04$. So about 90\% of excess activity variance is private to the agent. Talk shares more of its variance than activity ($\phi$ 0.21 against 0.08, in 77\% of 64 units), consistent with a read-out coupling on top of the scheduler. Practical value: trim each day to its all-present window, then report $c_\times$, before any claim about activity synchrony.

\begin{figure*}[t]
\centering
\includegraphics[width=\textwidth]{H38-platform-stalls/fig.pdf}
\caption{\textbf{The scheduler is a field on activity (H38).} (a)~In a simulated scheduled day, shared start and stop edges alone inflate the co-activation gain (0.42 over the whole day against 0.15 inside the all-present window). (b)~Village, round 1b: excess co-activation gain per period over the block-shift null (gray, 95\%), split into the scheduler part (orange) and the part that survives (blue); the median scheduler share is 0.72 in regime III and 0.11 in regime I.}
\label{fig:fields-scheduler}
\end{figure*}

\subsection{A goal is a field, and a kickoff is a quench}

A kickoff message names a target. On day 1, a period's content centroid sits closer to its own kickoff text than to the 32 other eligible kickoffs in 18 of 33 periods (median percentile 1.0, $p=3\times10^{-6}$; H54) \cred{0.92}. The misses are kickoffs that name no shared target: free weeks \#3 and \#16, \#44's half-free week and \#51. In \#51 each agent has a private goal text. An agent's content sits nearer its own goal than a swapped agent's goal in 0.95 of pairs ($p=0.0002$), with no decay over nine weeks.

The goal is not a small linear push. H10 tested the Legendre form $P_A\propto P_F\,e^{\lambda y}$, in which free-week fluctuations along the goal predict who moves. They do not. The combined test fails in all 8 input configurations (Stouffer $p$ 0.84--0.96), and the correlation is negative in both regime-I pairs ($r=-0.37$, $-0.69$; H10) \cred{0.79}. The kickoff push is 1.4--3.1 free-week standard deviations, far outside linear response. The assigned week fills content that the free week never visited.

The kickoff acts as a restoring force toward a new centre (Fig.~\ref{fig:fields-quench}). The cross-agent memory of pre-kickoff positions falls from $\rho_0=0.68$ across ordinary nights to $0.42$ at the kickoff. The extra forgetting is $\Delta\rho=+0.26\pm0.04$ (s.e.; 90\% CI $\pm0.07$), positive in 16 of 18 kickoffs, in both embedding models and all 7 input variants (H97) \cred{0.70}. The force is anisotropic: forgetting is $0.63\pm0.27$ along the kickoff direction and $0.25\pm0.04$ across it. Day 1 overshoots the settled level by $a=+0.12\pm0.02$. The per-agent susceptibility is not shown to be an agent constant (split-half $r=0.30$, $p=0.14$, power 0.30).

Seen from the old state, a goal switch is a quench, not a hysteresis loop. On day 1 the old state's field-orthogonal remanence is $R_1=0.27$ (bge) and 0.18 (gte) of its pre-switch level, against 0.85--0.89 across an ordinary night. Over 24 transitions, $\Delta R_1=-0.62$ [$-0.78$, $-0.45$] (bge; H96) \cred{0.74}. Most of the old state is gone within the first active half hour ($\tau_{\rm old}\approx0.15$ active h). Hysteresis predicts that a more ordered old state resists longer. That test has power 0.50, and its measured sign is reversed.

A trace of the previous period remains, but in the members. On day 1, agents load on the previous period's village centroid with an excess $\Delta\gamma_1=0.158$ (random effects 0.146 [0.080, 0.213], bge) over 25 boundaries. This is about three times the null 95th percentile of 0.055 (H82) \cred{0.73}. The trace does not decay over days 1--5. When the agent's own previous-period content enters the regression, the excess falls to $-0.02$ (post hoc). The remanence is each member's continuing themes, not the village record.

A failed hypothesis fixes the clock of this memory. H103 predicted that nights demagnetize content, as an alternating field demagnetizes a magnet. The night clock fits best in only 5 of 22 periods with decay (3 of 22 in gte). At fixed active lag, the night step is $\beta_N=-0.001$ [$-0.023$, $+0.021$] over 41 units (H103) \cred{0.59}. The synthetic clock-selection accuracy is 0.82, so this is a failure and not an inconclusive test. Content memory decays with active work, and an agent's own idle hour does not age it.

As a two-state system (on-goal or off-goal), a kickoff raises the on-goal occupancy from a decoy floor of about 0.03 to 0.23--0.44, always below one half, over 16 kickoffs. The variance grows as $p(1-p)$ (Spearman 0.86, $p=0.0003$; H105) \cred{0.36}. This is descriptive. Any binary occupancy gives it, and the tilt test that the model was meant to enable is not identifiable at village sampling. Practical value: an operator does not need to clear a swarm before a goal change, because the kickoff erases 70--80\% of the old state by day~1, most of it within the first active hour.

\begin{figure*}[t]
\centering
\includegraphics[width=\textwidth]{H97-quench-restoring-force/fig.pdf}
\caption{\textbf{A kickoff is an anisotropic restoring force (H97).} (a)~One kickoff (\#38$\to$\#39) in the plane of the new kickoff direction and the old-state direction, with ordinary-night moves in gray; (b)~cross-agent memory of positions falls from 0.68 across 161 ordinary nights (gray) to 0.42 at 18 kickoffs (black), $\Delta\rho=+0.26\pm0.04$. (c)~Extra forgetting is $0.63\pm0.27$ along the kickoff direction and $0.25\pm0.04$ across it, and day~1 overshoots the settled level by $a=+0.12\pm0.02$.}
\label{fig:fields-quench}
\end{figure*}

\subsection{Rooms are domains}

In regime III the agents sit in rooms (\#best and \#rest). A room could act as a field on its members or as a filter on what they read. For talk it is a filter. In a pooled panel of regime-III pair-days with pair and day fixed effects, co-location raises talk co-movement by $+0.0123$ ($z=5.1$). With ledger read counts in the model, the co-location term falls to $+0.0013$ ($z=0.2$), and reads carry the effect ($z=2.1$; H05) \cred{0.88}. Co-location and log reads correlate at 0.82, so the interval on the residual room term still includes half the room effect. Activity spins show no room effect after the trim.

Content coherence also stops at the room boundary. The ratio of cross-room to within-room per-pair content correlation is $C_B\le0.3$ with $p<0.05$ in 6 of 9 multi-room periods (median 0.23; H47) \cred{0.66}. At the NE42 merge and split, the boundary moves with the channel (difference in differences 1.0 [0.50, 1.34], $p=0.001$). Inside a room, pairs that address each other carry most of the coherence. Post hoc, the boundary is sharp only where rooms do different work. Human naming-game experiments show the same dependence on the channel: sparse networks freeze into local conventions, and well-mixed groups reach one convention~\cite{fd:centola2015}.

Rooms with identical kickoffs still separate in content (Fig.~\ref{fig:fields-rooms}). We remove each member's leave-period-out constant and the measured room-field directions. The \#best/\#rest separation still exceeds random regroupings ($Q_{\rm spont}>1$, $p<0.05$) in 5 of 7 regime-III periods (bge), and in 3 of 5 identical-kickoff periods in both models (H100) \cred{0.71}. Composition explains 0.03--0.18 of the separation where rooms separate. The split direction does not persist from one goal to the next: all six consecutive pairs have $|z|<2$, and $R(\#39,\#41)=-0.13$. Each goal breaks the symmetry again. Room-specific kickoffs about double the separation ($Q$ 6.85 and 5.09, against an identical-kickoff median of 1.68), but the split lies along the kickoff text only at chance. The residual is endogenous. It is not identified as coupling, and it can contain room-specific work fields that we do not measure.

The onset of the split has no single shape. In \#37 and \#39 it starts near zero and grows over the week (day-1 to final ratio $r_1=0.09$ and $-0.09$). In \#41 and \#42 it is 70--90\% formed on day 1, but only half of its day-1 direction is kept ($c_1$ 0.50--0.59; H107) \cred{0.35}. Within a period, the room direction persists from day to day (noise-corrected $P(1)$ 0.69--0.95 in 6 of 8 periods). A Goldstone mode lets an unfielded direction diffuse. The identical-kickoff periods do diffuse faster: $D_\theta=0.35$ per day against 0.075 with room kickoffs, a ratio $R_D=4.7$ [2.1, 5.9] (H108) \cred{0.42}. The contrast rests on \#38 (without it, $R_D=1.4$ in bge) and is confounded with split size, so we do not claim a Goldstone mode.

Where rooms do different work, content forms two domains. The bimodality of stayers on the inter-room axis is $D=5.3$--6.1 in \#38, \#41, \#44 and \#51g (relabel 95th percentile about 1), against a median of 1.70 in identical-kickoff weeks. The wall between the domains has no interior. Agents who hop between rooms do not sit at intermediate positions; their content follows the room they speak in. Hopping read-outs do not move a hopper's home-room content: $\kappa_R=-0.014$ [$-0.042$, $0.024$] per e-fold of read-outs in \#51g, which excludes $\kappa\ge0.1$ at synthetic power $\ge0.93$ (H102) \cred{0.51}. The original prediction, hoppers inside a wall of finite width, failed by its kill rule.

Room order is not stored in the context window. Over 3,477 forced erasures in 8 regime-III two-room periods, the drop in an agent's alignment with its room's spontaneous axis is $\delta_F=-0.08$ [$-0.23$, $0.07$] (bge; gte $-0.07$ [$-0.18$, $0.04$]). This excludes a 30\% drop at synthetic power 1.00 (H109) \cred{0.62}. The erasure scrambles one information channel~\cite{fd:kolchinsky2018}, and the room order survives it. The order must sit in the agent's work, prompt or memory.

Merging rooms does not add a coupling. In NE42's merged week the talk coupling between former \#best and \#rest agents did not rise (merge contrast $-0.02$ [$-0.13$, $0.09$] per pair; power 0.95 at $J=0.3$). The read-gated coupling of pairs that shared a room all along was 0.08 [$-0.16$, $0.31$] in the merged week and 0.47 [0.27, 0.66] after the split (H119) \cred{0.54}. A read path is necessary for coupling but not sufficient, because a larger room dilutes the per-pair coupling. H119's prediction, that the fitted adjacency follows the merge, failed by its kill rule.

Rooms forked from one state also separate fast. At NE15 the role-play game forked into one copy per room. In the first active hour the two rooms' content overlap is $q=0.51$, against 0.80 for identical replicas. After day 1 the rooms share no same-day state (memory 0.07 [$-0.03$, $0.17$]; H118) \cred{0.52}. Practical value: to keep two groups' content separate, give them different work, not only different rooms ($D$ 5--6 against 1.5--2.4).

\begin{figure*}[t]
\centering
\includegraphics[width=\textwidth]{H100-room-symmetry-breaking/fig.pdf}
\caption{\textbf{The room content split is spontaneous and regenerates at each goal (H100).} (a)~Two-block soft-spin simulation with and without a room field; (b)~relabel excess $Q$ per period before and after removal of composition and measured field directions (gray: null 95th percentile); (c)~composition and spontaneous shares; (d)~remanence of the split direction between consecutive periods against the joint-relabel null. After both removals the rooms differ beyond random groupings in 5 of 7 regime-III periods, and no split direction survives a goal change.}
\label{fig:fields-rooms}
\end{figure*}

\subsection{Random fields and style}

In \#51 each agent has a private goal, which acts as a quenched random field. After style removal, the disorder ratio is $R=0.67$--0.73 in all five \#51 units, against 0.24--0.53 in shared-goal weeks. Role text explains 5--14\% of it, and the mean-field gain is $b=0.23$--0.31, with a day-bootstrap CI above 0 in 5 of 5 units (H98) \cred{0.49}. The pull is weak and far from the mean-field instability. Without style removal the shared-goal weeks rise to 0.74--0.79, so their agent-level disorder is mostly style. The pre-registered test of independent agents in this field failed: day-to-day overlaps synchronize beyond the shift null in 2 of 5 units.

Model families look like fields in content, but the family field is writing style. It appears in 9 of 15 periods and survives style residualization in none, in both embedding models (H13) \cred{0.87}. In behavior the family trace is weak (random effects 0.075 [0.030, 0.121], significant in 3 of 15 units). Style itself is mostly an agent constant. Three components explain a median 0.72 (IQR 0.65--0.80) of systematic style variance in 34 of 34 periods. The agent constant takes 0.55, and context fill takes 0.08 (0.011 in regime III; H73) \cred{0.80}. Style is not conserved, though. It drifts as the context fills and returns toward the agent's mean at a forced erasure, where content does not move (boundary percentile 0.56 for style against 0.52 for content; style above one half in 18 of 23 units; H46) \cred{0.75}. Across a goal switch, style still identifies agents at 0.76, against 0.51 for content (chance 0.15).

A shared model prior is one of the four impostors. Populations of identical LLMs can develop a collective bias with no bias in a single agent's first choice~\cite{fd:ashery2025}. In the village the agents differ, and the per-agent style constant is the largest field term in style. Residualize on it before reading any family or convention effect.

\subsection{Fields drift: the stationary window}

The rule to fit models within a period assumes a stationary state inside it. That assumption fails for long periods. In \#38 (17 days), the kinetic-Ising parameters of 1-min activity and talk spins differ between halves about twice as much as between random day halves. The drift ratio is $R=1.95$ for activity and 2.18 for talk (Holm-corrected $p\le0.008$), and the \#51 main body drifts by the same amount ($R=2.17$, 2.14; H120) \cred{0.68}. The drift is gradual: the NE17 boundary ranks 0.36 among \#38's 14 day boundaries. The 8-day unit 38a shows no drift ($R=1.28$, $p=0.13$, power 0.85). Drift in entropy production is untestable at these counts with a held-out version of the maximum-entropy bound of Ref.~\cite{fd:aguilera2026} (power $\le0.30$).

So the stationary window is about one week. A fit on a longer unit puts a drifting field into the fitted couplings. Practical value: split a period longer than about eight days into units of eight days or less before any within-unit fit of $J$ or $\mathbf h$.
